\section{Illustrative application: a self-consistent estimate}\label{sec:application}

The self-consistent scheme needs all four operations of the algebra on a complete cubic
tensor. It yields three results. The concentration tensor of the grain comes in blocks,
and it shows that the scheme needs six conditions where five have been imposed
(section~\ref{sec:application-blocks}). For grains of physical size the scheme has a
closed limit, in which the polycrystal sees the grain only through its Voigt estimate
(section~\ref{sec:application-limit}). And the average of the Eshelby tensor on which the
scheme rests is obtained exactly (appendix~\ref{app:eshelby}).

Solyaev~\cite{Solyaev2023} extended Kr\"oner's self-consistent scheme to strain-gradient
elasticity. A spherical grain of radius $R$, with tensors $\tC_\gr$ and $\tD_\gr$, lies in
the isotropic effective medium $(\tC,\tD)$ under a uniform strain and a uniform strain
gradient, and the averages over the grain and over its orientations must return them. He
works in form~I with the variable $\tilde\kappa_{ijk}=u_{k,ij}=\eta_{kij}+\eta_{kji}-\eta_{ijk}$,
that is, $\tilde{\bm\kappa}=(3\tP{\mathcal S}-2\tI)\sgrad$, which multiplies the copies in $\mathcal S$ by $1$ and those
in $\mathcal R$ by $-2$: a form-I block is the form-II block rescaled copy by copy, and the
algebras do not change. In this form, the gradient condition reads~\cite[eqs.~(7)
and~(10)]{Solyaev2023}
\begin{equation}
\tB=\big(\tI+\tSav\,\tD^{-1}(\tD_\gr-\tD)\big)^{-1},\qquad \avg{\tB}=\tI,
\label{eq:sc-condition}
\end{equation}
with $\tSav$ the average over the sphere of the sixth-order Eshelby tensor of the medium:
as in Solyaev's scheme, the Eshelby field, which is not uniform in the grain, is replaced
by its average. The fourth-order analogue, with $\tC$ and $\tC_\gr$, becomes Kr\"oner's condition for
large grains.

\subsection{The concentration tensor in blocks}\label{sec:application-blocks}

Since $\tD$ and $\tSav$ are isotropic, $\mathbb L=\tSav\tD^{-1}$ is an isotropic element
without major symmetry, and $\tB$ lies in $\Alg{\clscub}$ for a cubic grain, with
\begin{equation}
\tB^p=\big(I+\mathbb L^p(\tD^p_\gr-\tD^p)\big)^{-1},\label{eq:sc-blocks}
\end{equation}
where, by~\eqref{eq:cub-iso-units}, an isotropic element $(m_3,m_2,M)$ has the cubic blocks
$\mathrm{diag}(m_3,M)$ in $\irr{T_1}$, $\mathrm{diag}(m_3,m_2)$ in $\irr{T_2}$, $m_2$ in
$\irr{E}$ and $m_3$ in $\irr{A_2}$. The grain enters through one $3\times3$, one $2\times2$
and two scalar inversions. By~\eqref{eq:weighted-means}, which holds without major
symmetry, $\avg{\tB}$ has the isotropic blocks
$y_3=\tfrac17(3B^{T_1}_{hh}+3B^{T_2}_{33}+B^{A_2})$,
$y_2=\tfrac15(3B^{T_2}_{22}+2B^{E})$ and $Y=B^{T_1}_{(sr)}$, and
\begin{equation}
\avg{\tB}=\tI\iff y_3=y_2=1,\quad Y=I_2:\label{eq:sc-six}
\end{equation}
six scalar equations, as many as the dimension of $\Alg{\clsiso}$. Solyaev imposes five
linear invariants~\cite[eq.~(13)]{Solyaev2023}. The sixth is $B_2'=B_{ijjkki}$, the
transpose of his $B_2=B_{iikkjj}$, and in his variables
\begin{equation}
B_2-B_2'=3\sqrt5\,(Y_{rs}-Y_{sr}),\label{eq:sc-sixth}
\end{equation}
so the five fix $\avg{\tB}$ only if $Y$ is symmetric, which $\tSav$ does not ensure (its
$\Hrm1$ block is not symmetric; appendix~\ref{app:eshelby}). We measure the mismatch by the
residual $\scres$ of his eq.~(22), printed without squares, with every term squared: the
root sum of squares of the
relative errors of $A_1$, $A_2$ (fourth order) and $B_1,\dots,B_5$, seven terms, or eight
with $B_2'$.

\subsection{Large grains}\label{sec:application-limit}

Solyaev's medium is the Helmholtz tensor $\ell^2\tC\ot\kron$ of form~II (his eq.~(19), read
as in~\cite[eq.~(4.3)]{Ma2018} and~\cite[eq.~(3.3)]{Gao2009}; read literally with his
$\tilde\kappa=\nabla\nabla\disp$, it would put $\tC$ on the derivative indices, which is not
the Helmholtz tensor, and such readings are tried in section~S6 of the supplementary
material), taken to form~I, with the bulk modulus of the
grain, so that $A_1=3$; the unknowns are $\nu$ and $\ell$. For this medium the exact average is
$\tSav=\Phi(R/\ell)\,\hat{\tS}(\lambda,\mu)$ (appendix~\ref{app:eshelby}), with
$\Phi$ the function of~\cite[eq.~(21)]{Solyaev2023}, $\Phi\approx-\ell/(2R)$ for $R\gg\ell$. Write $\tD=\ell^2\hat\tD$.
Since the orientation average commutes with composition by isotropic elements, to first
order in $\Phi$
\begin{equation}
\avg{\tB}-\tI=-\Phi\,\hat{\tS}\hat\tD^{-1}\big(\ell^{-2}\Vg{\tD_\gr}-\hat\tD\big),
\label{eq:sc-limit}
\end{equation}
where only the Voigt estimate of the grain enters. The $A$ conditions fix $\nu$ at
Kr\"oner's value $\nu_{\mathrm K}$~\cite[eq.~(3.7)]{Blaschke2017},
\cite[eq.~(16b)]{Kube2016}, and the $B$ terms give, to leading order, $\scres=|\Phi|\,|\ell^{-2}\bm\alpha-\bm\gamma|$,
with $\bm\alpha$ and $\bm\gamma$ the invariants $B_t$ relative to their values $B_t^0$ for $\tI$, of
$\hat{\tS}\hat\tD^{-1}\Vg{\tD_\gr}$ and $\hat{\tS}$. As $|\Phi|$ grows with $\ell$, the
minimum is at
\begin{equation}
\linf^{-2}=|\bm\gamma|/|\bm\alpha|,\qquad \scres/|\Phi|\;\to\;\big|\linf^{-2}\bm\alpha-\bm\gamma\big|.
\label{eq:sc-linf}
\end{equation}

\subsection{Results}\label{sec:application-results}

Table~\ref{tab:application} gives the limit for the positive definite data. It agrees
with the full minimisation at $R=1000$~\AA, and from $2R=2$~nm on the minimising $\ell$ is
within $\NumSCNearLimitPercent$\,\% of $\linf$ (supplementary material, section~S5). Three
readings follow. First, for grains of physical size the scheme sees the grain only through
$\Vg{\tD_\gr}$: no inverse of $\tD_\gr$ enters, and the block inversions
of~\eqref{eq:sc-blocks} matter only at sub-nanometre sizes. Second, the limit diagnoses the
scheme itself. The Helmholtz medium has equal block lengths~\eqref{eq:block-lengths}, while
those of $\Vg{\tD_\gr}$, with Kr\"oner's $\tC$, spread by $\NumAlVoigtKronerSpread$,
$\NumCuMEAMVoigtKronerSpread$ and $\NumCuEAMVoigtKronerSpread$. No single $\ell$ meets the
six conditions: the normalised residual stays at
$\scres/|\Phi|=\NumAlLimResidualSeven$--$\NumCuEAMLimResidualSeven$. The residual itself
vanishes with $|\Phi|$: with seven terms it is $\NumSCResidualTwoNmMin$--$\NumSCResidualTwoNmMax$
at $2R=2$~nm and $\NumSCResidualHundredNmMin$--$\NumSCResidualHundredNmMax$ at $2R=100$~nm, so
for large grains every $\ell$ of the order of 1~\AA\ nearly meets the conditions,
and $\ell$ is only weakly determined, as the minimiser of the coefficient of a vanishing
quantity. That minimiser $\linf$, a weighted compromise, exceeds $\lLP$ of the same Voigt estimate by $\NumAlLimOverLengthPercent$,
$\NumCuMEAMLimOverLengthPercent$ and $\NumCuEAMLimOverLengthPercent$\,\%. A one-length
medium cannot represent these polycrystals, and a medium with five gradient constants is
the next step, for which appendix~\ref{app:eshelby} already gives the general average.
Third, the sixth condition $B_2'=3$, which the seven-term solutions miss by up to
$\NumSCSixthSolutionMaxPercent$\,\%, moves $\linf$ by at most $\NumCuEAMLimShiftPercent$\,\%. Like the averages, the scheme
ignores the moduli induced by the heterogeneity of $\tC$~\cite{Malyi2021}, which at the
grain sizes where it reduces to $\Vg{\tD_\gr}$ are the dominant ones
(section~\ref{sec:averages-cubic}): it illustrates the algebra, not a homogenisation
theory.

With the grain constants of Solyaev's iron example, and grain and medium both written as the
form~II Helmholtz model, the scheme returns $\ell=\NumFeAllEllMin$--$\NumFeAllEllMax\,\ell_0$
and Kr\"oner's $\nu$ over the grain sizes of his figure~2, with any of the three averaged
tensors (the exact one, that of Ma et al.\ and his printed one); we could not reproduce the
published increase of $\ell$ (supplementary material, section~S6).

\begin{table}[pos=tbp]
\caption{Self-consistent estimate for large grains~\eqref{eq:sc-linf}: Poisson's ratio
$\nu_{\mathrm K}$ of Kr\"oner's classical estimate, with the bulk modulus of the grain; length
$\linf$ and normalised residual $\scres/|\Phi|$ with seven terms and, in brackets, eight;
spread $\ell_{\max}/\ell_{\min}$ of the block lengths~\eqref{eq:block-lengths} of the Voigt
estimate and its $\lLP$~\eqref{eq:lazar-po-length}, both relative to Kr\"oner's $\tC$.
Digits as in table~\ref{tab:averages}.}\label{tab:application}
\centering\small
\input{application}
\end{table}
